\begin{abstract}

Разработчики программного обеспечения создают языки программирования для
конкретных целей, и эффективное использование этих языков требует качественной
поддержки при редактировании кода. Представленный в 2016 году протокол языкового
сервера (Language Server Protocol) призван помочь разработчикам получить
расширенную поддержку в предпочитаемом ими текстовом редакторе (vim, neovim) или
IDE (VSCode, Eclipse IDE) путём разделения «языковой» и «редакторской» частей.

Поскольку современные IDE и некоторые текстовые редакторы поддерживают протокол
языкового сервера «из коробки», для получения ряда функций редактирования
разработчикам достаточно реализовать всего один программный компонент — так
называемый языковой сервер. По какой-то причине не для всех языков
программирования создаётся такая поддержка. Одним из таких языков является
Jasmin. Jasmin — это язык описания криптографических примитивов, и его основными
пользователями являются исследователи в области криптографии. В данной работе я
представляю реализацию языкового сервера для языка программирования Jasmin,
который поддерживает диагностику, переход к определению, инкрементальный
синтаксический анализ, подсветку символов, подсветку синтаксиса и отображение
типа символа при наведении курсора.

\end{abstract}